% Part: first-order-logic
% Chapter: natural-deduction
% Section: identity

\documentclass[../../../include/open-logic-section]{subfiles}

\begin{document}

\olfileid{fol}{ntd}{ide}

\olsection{\usetoken{S}{identity} ધરાવતી \usetoken{P}{derivation}}

!!{identity} ધરાવતી !!^{derivation}s માટે વધારાના અનુમાન-નિયમો જરૂરી છે.

\begin{defish}
\AxiomC{}
\RightLabel{\Intro{\eq}}
\UnaryInfC{$\eq[t][t]$}
\DisplayProof
\hfill
\begin{tabular}{r}
\AxiomC{$\eq[t_1][t_2]$}
\AxiomC{$!A(t_1)$}
\RightLabel{\Elim{\eq}}
\BinaryInfC{$!A(t_2)$}
\DisplayProof
\\[3ex]
\AxiomC{$\eq[t_1][t_2]$}
\AxiomC{$!A(t_2)$}
\RightLabel{\Elim{\eq}}
\BinaryInfC{$!A(t_1)$}
\DisplayProof
\end{tabular}
\end{defish}

ઉપરના નિયમોમાં $t$, $t_1$ અને $t_2$ બંધ પદો છે. \Intro{\eq} નિયમ
આપણને કોઈ ધારણા વિના, સ્વરૂપ $\eq[t][t]$નું કોઈ પણ તાદાત્મ્ય-વિધાન
સીધું !!{derive} કરવાની છૂટ આપે છે.

\begin{ex}
જો $s$ અને $t$ બંધ પદો હોય, તો $!A(s), \eq[s][t] \Proves !A(t)$:
\begin{prooftree}
\AxiomC{$\eq[s][t]$}
\AxiomC{$!A(s)$}
\RightLabel{$\Elim{\eq}$}
\BinaryInfC{$!A(t)$}
\end{prooftree}
આ અભિન્ન વસ્તુઓની પ્રતિસ્થાપનીયતાના સિદ્ધાંત તરીકે, અથવા લાઇબ્નિત્સના
નિયમ તરીકે, પરિચિત હોઈ શકે છે.
\end{ex}

\begin{prob}
$=$ સંમિત અને પરંપરિત બંને છે એ સાબિત કરો, એટલે કે
$\lforall[x][\lforall[y][(\eq[x][y] \lif \eq[y][x])]]$ અને
$\lforall[x][\lforall[y][\lforall[z]((\eq[x][y] \land \eq[y][z])
\lif \eq[x][z])]]$ની !!{derivation}s આપો.
\end{prob}

\begin{ex}
આપણે !!{sentence}
\begin{align*}
& \lforall[x][\lforall[y][((!A(x) \land !A(y)) \lif \eq[x][y])]]
\intertext{ને !!{sentence}}
& \lexists[x][\lforall[y][(!A(y) \lif \eq[y][x])]]
\end{align*}
માંથી !!{derive} કરીએ છીએ.
આપણે !!{derivation}ને પાછળથી વિકસાવીએ છીએ:
\begin{prooftree}
\AxiomC{$\lexists[x][\lforall[y][(!A(y) \lif \eq[y][x])]]
\quad \Discharge{!A(a) \land !A(b)}{1}$}
\DeduceC{$\eq[a][b]$}
\DischargeRule{\Intro{\lif}}{1}
\UnaryInfC{$((!A(a) \land !A(b)) \lif \eq[a][b])$}
\RightLabel{\Intro{\lforall}}
\UnaryInfC{$\lforall[y][((!A(a) \land !A(y)) \lif \eq[a][y])]$}
\RightLabel{\Intro{\lforall}}
\UnaryInfC{$\lforall[x][\lforall[y][((!A(x) \land !A(y)) \lif \eq[x][y])]]$}
\end{prooftree}
હવે મુખ્ય ધારણાનો ઉપયોગ કરવો પડશે: તે અસ્તિત્વપરિમાણિત !!{formula}
હોવાથી વચલું ફલિતવાક્ય $\eq[a][b]$ !!{derive} કરવા
\Elim{\lexists} વાપરીએ છીએ.
\begin{prooftree}
\AxiomC{$\lexists[x][\lforall[y][(!A(y) \lif \eq[y][x])]]$}
\AxiomC{$\Discharge{\lforall[y][(!A(y) \lif \eq[y][c]])}{2}$}
\noLine
\UnaryInfC{$\Discharge{!A(a) \land !A(b)}{1}$}
\DeduceC{$\eq[a][b]$}
\DischargeRule{\Elim{\lexists}}{2}
\BinaryInfC{$\eq[a][b]$}
\DischargeRule{\Intro{\lif}}{1}
\UnaryInfC{$((!A(a) \land !A(b)) \lif \eq[a][b])$}
\RightLabel{\Intro{\lforall}}
\UnaryInfC{$\lforall[y][((!A(a) \land !A(y)) \lif \eq[a][y])]$}
\RightLabel{\Intro{\lforall}}
\UnaryInfC{$\lforall[x][\lforall[y][((!A(x) \land !A(y)) \lif \eq[x][y])]]$}
\end{prooftree}
ઉપર જમણી બાજુની ઉપ-!!{derivation}માં તેની ધારણાઓ વાપરી
$\eq[a][c]$ અને $\eq[b][c]$ બતાવીને કામ પૂરું થાય છે. એ માટે બે જુદી
!!{derivation}s જોઈએ. $\eq[a][c]$ માટેની !!{derivation} આ છે:
\begin{prooftree}
\AxiomC{$\Discharge{\lforall[y][(!A(y) \lif \eq[y][c]])}{2}$}
\RightLabel{\Elim{\lforall}}
\UnaryInfC{$!A(a) \lif \eq[a][c]$}
\AxiomC{$\Discharge{!A(a) \land !A(b)}{1}$}
\RightLabel{\Elim{\land}}
\UnaryInfC{$!A(a)$}
\RightLabel{\Elim{\lif}}
\BinaryInfC{$\eq[a][c]$}
\end{prooftree}
$\eq[a][c]$ અને $\eq[b][c]$માંથી \Elim{\eq} વડે $\eq[a][b]$
!!{derive} કરીએ છીએ.
\end{ex}

\begin{prob}
નીચેનાં !!{formula}sની !!{derivation}s આપો:
\begin{enumerate}
\item $\lforall[x][\lforall[y][((\eq[x][y] \land !A(x)) \lif !A(y))]]$
\item $\lexists[x][!A(x)] \land \lforall[y][\lforall[z][((!A(y) \land
    !A(z)) \lif \eq[y][z])]] \lif \lexists[x][(!A(x) \land
  \lforall[y][(!A(y) \lif \eq[y][x])])]$
\end{enumerate}
\end{prob}

\end{document}
